和柏宇见完面回到家的这段时日，我经历了漫长的失眠。我常常在凌晨三四点惊醒，坐在床上，盯着沙发靠背上那块擦不掉的咖啡渍发呆。我陷入了一种半睡半醒的状态，被染过的靠背似乎正一点点向我压过来。我的脑子就像一个被疯狂驱动的轮子，有时是各种音乐声音，有时是理佳说的话，有时是别的什么。我越想将它拔出来，它就扎得越深。

“啊！我受不了了！”

我狠狠地用手砸自己的脑袋，“怎么才能将这些该死的念头清出去。”

二十多年来，我简直就像一只上了发条的滑稽木偶，全凭着一种自欺欺人的惯性在泥沼里爬行。升学，工作，甚至是在人前维持那种体面得滴水不漏的社交微笑。理佳的出现，不过是我在绝望中抓到的一根稻草。可是现在，在这漫长得仿佛没有尽头、黏稠得让人窒息的黑夜里，幻象破灭了！我的脑子里空空荡荡，只觉得整个生命就像这深不见底的黑夜，既冗长，又荒谬得叫人发疯。

“亮妈妈，亮在语文课上看小说，被课任老师发现把书没收了。” 班主任从抽屉里拿出来一本书，交给了我的母亲。

“其实我们也不反对孩子看小说，这样可以培养语感......，只是在这个时间段，看小说有点耽误学习，而且我听说亮上课时候经常和旁边同学搭话，这也有点耽误其他孩子学习呀......”

我看到妈妈脸上露出尴尬的笑容，不断地向老师陪着笑脸。在回家的路上的时候，我看到妈妈一直阴沉着脸，我坐在公交车的靠里座，双腿蜷缩地不敢动弹，我眼睛瞥向妈妈，她面容严肃地像一个法官，不知道什么时候会审判我的罪行。说句话，哪怕骂我也好呢？我心里想着，可是我又做错了什么呢？

五年后，当我和爸爸妈妈说了我的志向后，他们用一种匪夷所思的目光看向我。

“要退学，你知道你在说些什么吗？你要写小说也没人管你，可你为什么要退学，你知道你现在的教育机会有多么来之不易吗。在我们那个年代，我们想上学都没得上，早早就出来工作。现在倒好，你有这么优越的条件，却偏偏不懂得珍惜，你现在退学了，以后就等着扫大街吧！”

我陷入了长久的失眠中，脑海中的念头犹如一首不停循环播放的音乐一般，但是我既没有暂停键也无法调控音量。在经历了几周难忍的日子后，我终于无法忍受，于是央求爸爸带我去看医生。

听完我的症状后，医生冷冰冰地看着我，用一种不屑一顾的口吻说道，“你这没什么症状，纯粹是想的太多了。”

突然，我有一种不可遏制的愤怒涌上心头，我说，“你配做个医生吗，你到底懂不懂什么是医学，你到底理不理解我的感受！”

某次，我坐在工位上，终于无法控制地大口呼吸起来，一道无法察觉的眼泪从我面颊划过，同事发现了我的异常，我的上司找我谈话，委婉地劝我休息一下。
